\documentclass[10pt]{article}
\usepackage[margin=20mm]{geometry}
\usepackage{amsmath,amssymb,fontspec,unicode-math,xurl,hyperref}
\IfFontExistsTF{DejaVu Serif}{\setmainfont{DejaVu Serif}}{\setmainfont{Latin Modern Roman}}
\setmathfont{latinmodern-math.otf}
\setlength{\parindent}{0pt}\setlength{\parskip}{6pt}\sloppy
\begin{document}
ELVIA Scientific Verification\par
I25-40103 — No Issue Found within inspected scope\par
Position: We Need Responsible, Application-Driven (RAD) AI Research\par
Review date: 23 October 2026 (America/Chicago). Exact assigned PMLR version; printed/PDF pages 1-12. AI-assisted review; scientific decision pending human adjudication.\par
\# 1. Identity, coverage and outcome\par
Reviewed: Sarah Hartman, Cheng Soon Ong, Julia Powles and Petra Kuhnert, PMLR 267:81514-81525 (2025), 12 pages, 2,389,590 bytes. SHA256: C2E8F1E2E88741F7A368C4DE62893193CF0D74A1F612DED77F06706AC3F5B5D7. Printed page numbers equal one-based PDF positions. The main argument and Figure 1, pp.1-9, were inspected; references on pp.9-12 were checked selectively. The assigned PDF contains no appendices.\par
No substantive mathematical, statistical, methodological or implementation defect was demonstrated within this scope. RAD-AI is an organizational research proposal, presented as a starting point in Section 3, p.5. The article supplies no new fitted statistical model, executable algorithm or controlled evaluation of the full framework. No author code or empirical reproduction was run, and no correction is proposed.\par
\# 2. Checked reasoning\par
Claim or component | Check and limit\par
Sections 1-2, pp.1-5: motivation | The synthesis of responsible research/innovation and application-driven AI is a normative position. Its necessity and expected societal benefits are not demonstrated universal theorems.\par
Stage 1, pp.5-6 | Team formation and people-centred design are coherent with the contextual aim. Human-centred design and community-based participatory research are distinguished and may be combined.\par
Stage 2, pp.6-7 | Methods, ethics, hidden assumptions and application metrics are considered together. White-box-only restrictions are rejected; concept leakage, memory use and vanishing gradients are acknowledged. Whether AI is suitable is explicitly questioned.\par
Stage 3, pp.7-8 | Testbeds and communities of practice are proposed for continued reassessment across time, scale and context. A successful testbed is not a proof of safe or beneficial performance in every deployment.\par
Section 4, p.8; Figure 1, p.5 | The critiques concern treating AGI, openness or a one-time checklist as sufficient. The paper retains useful roles for open development and structured tools; the figure agrees with the three-stage text.\par
Conclusion, p.9 | Claims of better trust, impact and research support are anticipated benefits. Their magnitudes, costs and causal effects remain unmeasured here.\par
\# 3. Selected primary-source checks\par
Agricultural emulation (Section 2.2, p.3): Powell et al.’s MODSIM 2023 abstract, p.699, describes a proof of concept using APSIM simulations for wheat in Dalby to train a neural biomass emulator with uncertainty estimates [2]. It supports the existence and intended role of model emulation. That one-page abstract provides no numerical runtime benchmark, so it alone does not quantify the speedup language. No broader claim about all cited emulators is inferred.\par
Broadcasting (Section 2.4, p.5; Section 3.2): ABC Align, arXiv:2408.00307v1, explicitly conditions work on organizational principles, content and editorial feedback [3]. Its Sections 5 and 9.1 distinguish standard benchmark evaluation from future organization-specific evaluations. It supports a concrete contextual-alignment example, not proof that the entire RAD workflow or every editorial obligation has been validated. Its model/data experiments were not reproduced.\par
Historical legislation (Section 2.3, p.3): AI Index 2024, Chapter 7, Section 7.2, chapter PDF p.11, reports that 32 of the 128 countries studied enacted at least one AI-related bill during 2016-2023 [4]. This supports the target’s over-30 statement by 2024. It counts enacted legislation mentioning AI; it does not imply 32 comprehensive AI regulatory regimes. This is a historical source check, not a current legal-compliance assessment.\par
The official ISO record confirms the cited ISO/IEC 42001:2023 management-system standard [5]; the full standard was not inspected. The AAAS resource was corroborated at the official search-index level, but its full page was inaccessible. The precise UNESCO consultation-paper p.45 testbed quotation was not independently verified. These access limits do not establish a misquotation or invalidate the workflow.\par
\# 4. Versions, disposition and remaining work\par
The official arXiv history lists v1 (7 May 2025), v2 (9 June 2025) and v3 (19 August 2025) [6]. The current metadata describes a final camera-ready version with updated figure color and reference. Those are author-provided change descriptions; no complete version diff or byte identity with PMLR is claimed. The assigned PMLR record links the PDF and OpenReview but no separate software or supplement. No separately declared substantive erratum was identified in the inspected records.\par
Item | Disposition\par
Scientific outcome | No Issue Found within inspected scope; not a claim of whole-paper correctness\par
Repair and author adoption | No correction proposed; adoption not requested\par
Remaining research | Compare concrete RAD implementations, costs and contextual outcomes; establish decision authority and test real deployment failures\par
Unperformed checks | Complete reference/standards audit, full version diff, cited experiments and author-code execution\par
Scientific decision | Pending human adjudication\par
The coverage, findings and provenance JSON record locations, sources, access limits, commands and hashes. Source snapshots and the visual record are in sources/0005. Two assistant reviewers checked conceptual reasoning and primary evidence; root reconciled the findings and produced the artifacts. No human-review declaration was completed. TeX and PDF were generated from common structured content; the PDF was rendered through the available HTML renderer, and TeX was not compiled. No software was installed.\par
\# 5. Consulted references\par
[1] Assigned PMLR proceedings record and PDF. https://proceedings.mlr.press/v267/hartman25a.html\par
Powell et al. (2023), Towards an AI Agronomist, MODSIM, p.699. https://mssanz.org.au/modsim2023/files/powell.pdf\par
Seneque et al., ABC Align, arXiv:2408.00307v1 (1 August 2024); selected methods/evaluation/limitations. https://arxiv.org/html/2408.00307v1\par
Stanford AI Index 2024, Chapter 7, Section 7.2, p.11. https://hai.stanford.edu/assets/files/hai\_ai-index-report-2024\_chapter\_7.pdf\par
ISO/IEC 42001:2023, official overview and metadata only. https://www.iso.org/standard/42001\par
Target preprint, arXiv:2505.04104, version history through v3. https://arxiv.org/abs/2505.04104\par
\end{document}
